\documentclass[11pt]{article}
\usepackage[a4paper,margin=1in]{geometry}
\usepackage[T1]{fontenc}
\usepackage{lmodern,microtype}
\usepackage{amsmath,amssymb,amsthm,mathtools}
\usepackage{booktabs,enumitem}
\usepackage{xcolor}
\usepackage[numbers,sort&compress]{natbib}
\usepackage[colorlinks=true,linkcolor=blue!55!black,citecolor=blue!55!black]{hyperref}

\newtheorem{theorem}{Theorem}[section]
\newtheorem{corollary}[theorem]{Corollary}
\theoremstyle{remark}
\newtheorem{remark}[theorem]{Remark}
\newcommand{\norm}[1]{\left\lVert#1\right\rVert}

\title{Why Radial Shell Gaps Do Not Control Nonnormal Riesz Projections}
\author{Liang Wang}
\date{July 2026}

\begin{document}
\maketitle

\begin{abstract}
The shell-complete resonance atlas has positive radial gaps, yet its
biorthogonal cloud projectors become extremely ill-conditioned.  We prove
that there is no contradiction.  For
\[
 A_M=\begin{pmatrix}\lambda&M\\0&\mu\end{pmatrix},\qquad \lambda\ne\mu,
\]
the spectral projector onto $\lambda$ is
\[
 P_{\lambda,M}=\begin{pmatrix}1&M/(\lambda-\mu)\\0&0\end{pmatrix},
\]
so $\norm{P_{\lambda,M}}_2=\sqrt{1+|M/(\lambda-\mu)|^2}$.  The eigenvalue gap
is fixed while the projector norm is unbounded.

In the RH-232 audit the minimum radial shell gap remains
$7.40\times10^{-5}$, whereas the maximum projector norm reaches
$2.26\times10^{12}$ and the minimum gap-to-projector ratio is
$6.40\times10^{-16}$.  Log--log fits against the noise give descriptive
growth exponents $7.41$ and $6.55$ in the two channels.  Radial separation is
therefore not a pseudospectral certificate.  The result marks a forbidden
shortcut but leaves projection-free determinant factorization open.
\end{abstract}

\section{Spectral separation versus invariant-subspace separation}

For normal operators a spectral gap controls the orthogonal spectral
projection.  For nonnormal operators, eigenvalue location and invariant-space
angle are different data \citep{Kato1995,TrefethenEmbree2005}.  RH-222
recorded the first; RH-232 measured the second
\citep{WangRH222,WangRH232}.

The simplest triangular family already contains the entire obstruction.

\section{Exact fixed-gap counterexample}

\begin{theorem}[Fixed-gap projector blow-up]\label{thm:fixed-gap}
Let $\lambda\ne\mu$ and
\[
 A_M=\begin{pmatrix}\lambda&M\\0&\mu\end{pmatrix}.
\]
The Riesz projector associated with $\lambda$ is
\[
 P_{\lambda,M}=\frac{A_M-\mu I}{\lambda-\mu}
 =\begin{pmatrix}1&M/(\lambda-\mu)\\0&0\end{pmatrix}.
\]
Consequently
\[
 \norm{P_{\lambda,M}}_2
 =\sqrt{1+\frac{|M|^2}{|\lambda-\mu|^2}}
 \longrightarrow\infty
\]
as $|M|\to\infty$, while the eigenvalue gap $|\lambda-\mu|$ is unchanged.
\end{theorem}

\begin{proof}
The polynomial formula follows because the minimal polynomial is
$(z-\lambda)(z-\mu)$.  The displayed projector has only one nonzero row, so
its operator norm is the Euclidean norm of that row.
\end{proof}

\begin{corollary}\label{cor:radial}
No bound depending only on a positive eigenvalue gap, radial or otherwise,
can control the norm of a nonnormal spectral projector.
\end{corollary}

This statement does not say that gaps are useless.  They identify a contour
that avoids the spectrum.  What they do not control is the resolvent on that
contour, and the Riesz formula integrates the resolvent rather than the set of
eigenvalues alone.

\section{The missing datum is a resolvent bound}

\begin{theorem}[Contour-resolvent projector bound]\label{thm:contour}
Let $\Gamma$ be a positively oriented rectifiable contour in the resolvent
set of a bounded operator $A$, enclosing an isolated spectral component.
Its Riesz projector satisfies
\[
 P_\Gamma=\frac{1}{2\pi i}\int_\Gamma(zI-A)^{-1}\,dz,
 \qquad
 \norm{P_\Gamma}
 \le \frac{\operatorname{length}(\Gamma)}{2\pi}
       \sup_{z\in\Gamma}\norm{(zI-A)^{-1}}.
\]
Consequently, a uniform contour length and a uniform resolvent bound do give
a uniform projector bound; a spectral gap alone supplies neither.
\end{theorem}

\begin{proof}
The first identity is the Riesz functional calculus.  Taking norms under the
Bochner integral gives the second inequality.
\end{proof}

For the triangular family the resolvent can be written explicitly:
\[
 (zI-A_M)^{-1}
 =\begin{pmatrix}
 (z-\lambda)^{-1} & M((z-\lambda)(z-\mu))^{-1}\\
 0 & (z-\mu)^{-1}
 \end{pmatrix}.
\]
On the circle $|z-\lambda|=r<|\lambda-\mu|$, its norm is at least the
modulus of the upper-right entry, and hence at some---indeed every---point
obeys
\[
 \norm{(zI-A_M)^{-1}}
 \ge \frac{|M|}{r(|\lambda-\mu|+r)}.
\]
The contour remains a fixed positive distance from both eigenvalues while
the resolvent diverges linearly in $|M|$.  The same off-diagonal coupling that
creates the oblique projector therefore appears as pseudospectral inflation.

\begin{remark}
If $A=V\Lambda V^{-1}$ is diagonalizable, the elementary estimate
\[
 \norm{(zI-A)^{-1}}
 \le \frac{\norm V\,\norm{V^{-1}}}
          {\operatorname{dist}(z,\sigma(A))}
\]
shows exactly which datum supplements the spectral gap: a controlled
eigenbasis condition number.  In infinite-dimensional transfer problems an
adapted resolvent estimate is the more invariant formulation.
\end{remark}

\section{Archived dynamical audit}

The finite model uses $\lambda=0.8$, $\mu=0.5$, and couplings from $1$ to
$10^9$.  The gap stays $0.3$ while the projector norm grows by a factor
$9.58\times10^8$.

For the RH-232 endpoint family, the projector norms grow much faster than the
radial gaps shrink.  Representative values are shown below.

\begin{table}[ht]
\centering
\small
\begin{tabular}{rrrr}
\toprule
$\sigma$ & left $\norm P_2$ & right $\norm P_2$ & conclusion\\
\midrule
$0.04$ & $4.16\times10^1$ & $4.07\times10^1$ & already oblique\\
$0.008$ & $2.79\times10^5$ & $1.10\times10^6$ & unstable\\
$0.0032$ & $4.34\times10^9$ & $3.62\times10^9$ & severe\\
$0.00125$ & $2.26\times10^{12}$ & $5.44\times10^{10}$ & extreme\\
\bottomrule
\end{tabular}
\caption{Projector norms inherited from RH-232.}
\end{table}

The fitted power exponents are diagnostics, not asymptotic theorems: the log
residuals are large and the rank changes with the noise.  The robust
conclusion is only that no bounded trend is visible and radial gaps alone
cannot supply one.

\section{Scope of the negative result}

Corollary~\ref{cor:radial} is a no-go theorem for any argument whose only
input is eigenvalue location.  It does not prohibit estimates that use
positivity, reversibility, a symmetrizing density, analytic distortion, or a
semiclassical anisotropic norm.  Nor does it prove that the archived
projectors actually diverge in the continuum limit: the finite exponents are
descriptive and the matrices change dimension with $\sigma$.

The practical distinction is useful.  A future positive result must expose
where its resolvent control enters; it cannot cite the shell gap as a proxy.
Conversely, failure of the Euclidean projector does not contaminate spectral
quantities, such as a canonical determinant product, that depend only on
algebraic eigenvalues.

\section{Route consequence}

The RH-222 shell gaps remain valid finite spectral observations.  What fails
is the implication
\[
 \text{positive radial gap}
 \quad\Longrightarrow\quad
 \text{uniformly bounded Riesz projector}.
\]
Any operator-level continuation must add a resolvent estimate, an adapted
norm, or a different factorization principle.  RH-234 uses the canonical
eigenvalue product for $\det_2$, which evaluates a finite cloud factor without
ever multiplying by the ill-conditioned projector.

This paper does not exclude a certified contour on a different Banach space
and does not close Gate A or any later arithmetic gate.

\bibliographystyle{plainnat}
\bibliography{references}
\end{document}
